% !TEX root = ../main.tex
\section{배경}
\label{sec:background}

\subsection{Stacks와 Clarity}
\label{sec:bg-clarity}

Stacks는 비트코인에 정산하는 스마트 컨트랙트 블록체인이다 \cite{ali2021stacks}. 컨트랙트는 Clarity로 작성한다. Clarity에는 재귀와 무한 반복이 없어서 호출 비용의 상한을 실행 전에 정할 수 있다 \cite{sip002}. Clarity는 컴파일하지 않는다. 노드는 소스 텍스트를 저장하고 해석하며, Stacks Blockchain API는 배포된 어느 컨트랙트에 대해서든 그 텍스트와 컨트랙트 인터페이스(ABI)를 돌려준다 \cite{hiroapi}.

컨트랙트는 \code{SP\ldots ABC.amm-pool} 꼴의 컨트랙트 principal로 식별한다. 점 앞의 주소는 그 컨트랙트를 배포한 standard principal이므로, 식별자 자체에 코드를 게시한 사람이 기록된다. 컨트랙트는 \code{contract-call?}로 다른 컨트랙트를 호출한다. \emph{정적 디스패치}에서는 피호출자를 principal 리터럴로 쓰고, \emph{동적 디스패치}에서는 피호출자가 trait 타입의 함수 인자여서 호출하는 쪽이 실행 시점에 구체적인 컨트랙트를 고른다 \cite{sip002}. trait은 \code{define-trait}로 선언한다. 컨트랙트는 \code{impl-trait}로 어떤 trait을 따른다고 밝히고, \code{use-trait}로 trait을 가져와 인자 타입으로 쓴다. 두 구문 모두 그 trait을 정의한 컨트랙트를 가리킨다. Clarity~5는 상수도 \code{contract-call?}의 대상으로 받는다 \cite{sip039}. 따라서 정적 디스패치는 소스에서 읽어 낼 수 있는 호출 그래프를 만들고, 동적 디스패치는 트랜잭션이 피호출자를 넘겨줄 때까지 그 자리를 비워 둔다.

토큰 표준은 trait이다. SIP-009는 대체 불가능 토큰을 \cite{sip009}, SIP-010은 대체 가능 토큰을 \cite{sip010}, SIP-013은 준대체 가능 토큰을 정의한다 \cite{sip013}. Stacks Blockchain API는 ABI가 주어진 trait과 맞는 컨트랙트를 나열할 수 있어서 \cite{hiroapi}, 레지스트리는 컨트랙트를 파싱하지 않고도 표준을 따르는 컨트랙트에 표시를 붙일 수 있다.

\subsection{Stacks가 자산 이동을 승인하는 방식}
\label{sec:bg-authorization}

\textbf{SIP-010에는 allowance가 없다.} SIP-010 trait은 금액, 보내는 쪽, 받는 쪽, 선택적 메모를 받는 \code{transfer} 함수와, 이름·심볼·소수 자릿수·잔액·공급량·토큰 URI를 읽는 함수들을 선언한다 \cite{sip010}. 한 principal이 다른 principal의 토큰을 쓰게 하는 함수는 없다. 흔한 구현은 보내는 쪽이 트랜잭션 발신자일 때만 전송을 받아들인다. Clarity에서 \code{tx-sender}는 컨트랙트가 맥락을 자기 자신으로 바꾸지 않는 한 중첩 호출을 거쳐도 트랜잭션에 서명한 계정으로 남고, \code{contract-caller}는 바로 앞의 호출자이다 \cite{sip002}. 그래서 사용자가 호출한 탈중앙 거래소는 미리 승인을 받지 않고도 그 호출 안에서 사용자의 토큰을 옮길 수 있다. 사용자를 지키는 것은 포스트컨디션 메커니즘이다.

\textbf{포스트컨디션.} Stacks 트랜잭션은 포스트컨디션 목록과 포스트컨디션 모드를 담고, 둘 다 서명 대상에 포함된다 \cite{sip005}. 각 포스트컨디션은 principal(트랜잭션 발신자, standard principal, 컨트랙트 principal 중 하나), 자산(STX, 대체 가능 토큰, 특정 대체 불가능 토큰 중 하나), 조건을 지정한다. STX와 대체 가능 토큰의 조건은 그 principal이 트랜잭션 동안 보내는 양을 정해 둔 양과 $=$, $>$, $\geq$, $<$, $\leq$로 비교한다. 대체 불가능 토큰의 조건은 그 principal이 해당 토큰을 보낼지 보내지 않을지를 정한다. 실행이 끝나면 노드가 모든 포스트컨디션을 검사한다. 하나라도 어긋나면 트랜잭션은 중단되고 그 효과는 되돌려진다. 다만 트랜잭션은 블록에 남고 수수료는 낸다. 모드는 어느 포스트컨디션에도 나오지 않는 전송을 어떻게 다룰지를 정한다. \emph{Allow} 모드에서는 그런 전송을 허용하고, \emph{Deny} 모드에서는 트랜잭션을 중단한다 \cite{sip005}. Stacks.js는 따로 지정하지 않으면 Deny 모드로 트랜잭션을 만든다 \cite{hiropostconditions}. Clarity~5와 함께 Epoch~3.4(비트코인 블록 943,333, 2026년 4월)에 적용된 SIP-040은 세 번째 모드인 \emph{Originator}를 더한다. 이 모드는 트랜잭션 발신자의 자산에는 Deny 규칙을, 다른 모든 principal에는 Allow 규칙을 적용한다. SIP-040은 대체 불가능 토큰에 \emph{may send} 조건도 더했다 \cite{sip040,sip039}. 도입 동기는 DeFi 호출이 사용자가 미리 알 수 없는 방식으로 컨트랙트 사이에 자산을 옮겨서 사용자를 Allow 모드로 내몰았다는 데 있다.

\textbf{컨트랙트의 자산 한도.} 2025년 11월 비트코인 블록 923,222에서 활성화된 Clarity~4(SIP-033)는 컨트랙트가 자기가 통제하지 않는 코드를 실행하는 동안 자산 유출을 제한할 수 있게 한다 \cite{sip033}. \code{restrict-assets?}는 principal 하나, 한도 목록, 본문을 받아 본문을 실행한 뒤 그 principal의 유출을 한도와 비교한다. \code{as-contract?}는 \code{tx-sender}와 \code{contract-caller}를 컨트랙트 자신으로 바꾸어 본문을 실행하고, 컨트랙트 자신의 유출을 검사한다. 한도 표현식은 \code{with-stx}, \code{with-ft}, \code{with-nft}, \code{with-stacking}이며 각각 자산 하나와 그 한도를 지정한다. \code{as-contract?}에서만 쓸 수 있는 \code{with-all-assets-unsafe}는 모든 한도를 없앤다. 2026년 7월 비트코인 블록 960,230에서 Epoch~4.0과 함께 적용된 Clarity~6은 \code{with-stacking}의 이름을 \code{with-staking}으로 바꾸고, PoX 컨트랙트 호출을 위한 \code{with-pox}를 더했다. 새로운 종류의 트랜잭션 포스트컨디션은 더하지 않았다 \cite{sip044}. 유출이 한도를 넘으면 본문의 효과는 되돌려지고, 표현식은 어긴 한도의 순번을 오류로 돌려준다. 한도가 없는 자산이 움직이면 \code{(err u128)}을 돌려준다. Clarity~4는 아무 한도 없이 맥락만 바꾸던 옛 \code{as-contract}를 없앴다. 코드~\ref{lst:as-contract}가 이 패턴을 보여 준다. SIP-033은 \code{contract-hash?}도 더했다. 이 함수는 컨트랙트 코드 본문의 SHA-512/256 해시를 돌려주어, 한 컨트랙트가 다른 컨트랙트가 알려진 템플릿을 따르는지 확인할 수 있게 한다. 키워드 \code{current-contract}도 이때 더해졌다.

\begin{lstlisting}[language=Clarity,caption={Clarity~4의 자산 한도. 볼트 소유자만 지급을 시작할 수 있고, 본문이 실행되는 동안 볼트에서 나갈 수 있는 토큰은 최대 \code{amount}이다.},label={lst:as-contract},float=htbp]
(define-constant OWNER tx-sender)
(define-public (pay-out (amount uint) (recipient principal))
  (begin
    (asserts! (is-eq contract-caller OWNER) (err u401))
    (as-contract? ((with-ft .token "tok" amount))
      (try! (contract-call? .token transfer amount current-contract recipient none)))))
\end{lstlisting}

\textbf{컨트랙트 사이의 지속적인 승인.} 일부 프로토콜은 신뢰하는 컨트랙트의 목록을 명시적으로 둔다. 여러 Stacks DAO가 쓰는 ExecutorDAO 패턴에서는 핵심 컨트랙트가 활성화된 확장의 맵을 두고, 거버넌스가 \code{set-extension}을 호출해 확장을 켜거나 끈다. 켜진 확장은 \code{contract-caller} 검사를 거쳐 DAO의 송신 맥락에서 행동할 수 있다 \cite{executordao}. 이런 맵은 Stacks에서 지속적인 allowance에 가장 가까운 것이다. 명시적인 트랜잭션이 설정하고, 바뀔 때까지 효력이 남는다.

\subsection{패키지 레지스트리와 인기 신호}
\label{sec:bg-registries}

패키지 레지스트리는 게시된 코드를 색인하고 텍스트 관련도와 인기를 섞어 검색 결과의 순위를 매긴다. Go 패키지 사이트는 질의와 맞는 모듈마다 다음 점수를 계산한다.
\begin{equation}
\label{eq:pkgsite}
\mathrm{score} = \mathrm{ts\_rank}(q,m)\cdot\bigl(\ln(e + \mathrm{imported\_by}(m))\bigr)^{\alpha}\cdot\rho_{\mathrm{redist}}(m)\cdot\rho_{\mathrm{mod}}(m)
\end{equation}
여기서 $\mathrm{ts\_rank}$는 PostgreSQL의 전문 검색 순위, $\mathrm{imported\_by}$는 그 모듈을 가져오는 패키지 수이고, 마지막 두 인자는 재배포할 수 없거나 \code{go.mod} 파일이 없는 모듈에 벌점을 준다 \cite{pkgsite}. 로그 때문에 가져오는 패키지 수가 두 배가 되면 그 인자는 두 배가 되지 않고, 그 로그 값이 약 $\ln 2$만큼 커질 뿐이다. npm, Cargo, CRAN 등 여러 생태계의 의존성 네트워크는 그래프로 연구되어 왔고 \cite{kikas2017structure,decan2019empirical}, npm 의존성 그래프 위의 PageRank로 쇠퇴하는 패키지를 찾은 연구도 있다 \cite{mujahid2021centrality}.

tea 프로토콜은 의존성 그래프에서 패키지가 차지하는 자리를 점수로 매긴 teaRank에 따라 보상을 지급했다 \cite{tea2022}. \ref{sec:introduction}절에서 본 사건들은 이런 점수에 돈이 걸리면 무슨 일이 생기는지 보여 준다. 서로를 의존하는 스팸 패키지가 쏟아진다 \cite{sonatype2024tea,aws2025tea}. 레지스트리는 타이포스쿼팅과 계정 탈취를 비롯한 다른 공급망 공격도 겪으며 \cite{ohm2020backstabber,ladisa2023sok,vu2020typosquatting}, 소수의 관리자 계정이 전이 의존성을 거쳐 npm의 큰 부분에 닿을 수 있다 \cite{zimmermann2019small}.

\subsection{PageRank, 재시작 분포, 시빌 공격}
\label{sec:bg-pagerank}

$G=(V,E,w)$를 간선 가중치가 양수인 방향 그래프라 하고, $P$를 나가는 간선이 있는 노드에서는 $P_{uv}=w(u,v)/\sum_x w(u,x)$이고 나가는 간선이 없는 노드에서는 행 전체가 0인 행렬이라 하자. 감쇠 계수 $d\in(0,1)$와 $V$ 위의 재시작 분포 $\boldsymbol{\pi}$가 주어지면 PageRank는 다음 식의 고정점이다.
\begin{equation}
\label{eq:pagerank}
\mathbf{r} = (1-d)\,\boldsymbol{\pi} + d\,P^{\top}\mathbf{r} + d\,\delta(\mathbf{r})\,\boldsymbol{\pi},
\qquad
\delta(\mathbf{r}) = \sum_{v:\,\text{나가는 간선 없음}} r_v
\end{equation}
즉 나가는 간선이 없는 노드(댕글링 노드)가 가진 질량은 재시작 분포로 돌아간다 \cite{page1998pagerank,langville2006google}. 원래 알고리즘은 균등한 $\boldsymbol{\pi}$를 쓴다. 주제 민감 PageRank와 개인화 PageRank는 균등하지 않은 $\boldsymbol{\pi}$를 쓰고 \cite{haveliwala2003topic,jeh2003scaling}, TrustRank는 사람이 평판이 좋다고 판단한 소수의 페이지에 재시작을 몰아준다 \cite{gyongyi2004trustrank}. EigenTrust는 미리 신뢰한 피어에 같은 역할을 맡긴다 \cite{kamvar2003eigentrust}. AWP와 EndorseRank는 $d=0.85$와 균등 재시작을 쓴다 \cite{do2023awp,kwon2026endorserank}.

시빌 공격은 한 주체가 통제하는 신원을 많이 만드는 공격이다 \cite{douceur2002sybil}. PageRank를 상대로 하면 링크 팜의 꼴이 된다. 링크 팜은 대상을 가리키려고만 존재하는 페이지 집합이며, 여러 팜이 손을 잡으면 참여한 모든 팜의 점수가 오른다 \cite{gyongyi2005alliances}. Cheng과 Friedman은 노드 신원에 대해 대칭인 평판 함수는 시빌 공격에 안전할 수 없고, 신뢰하는 출발점에서 계산하는 비대칭 흐름 기반 함수는 안전할 수 있음을 보였다 \cite{cheng2005sybilproof}. SybilGuard와 SybilRank 같은 소셜 그래프 방어도 같은 관찰에 기댄다. 공격자는 노드와 그 사이의 간선을 마음대로 만들 수 있지만, 정직한 노드에서 공격자 영역으로 들어가는 \emph{공격 간선}의 수는 제한되어 있고, 정직한 노드에서 출발한 걷기는 이 간선을 좀처럼 건너지 않는다 \cite{yu2006sybilguard,cao2012sybilrank}. 커뮤니티 단위로 PageRank를 분석한 연구는 점수가 들어오는 링크에 어떻게 의존하는지를 정확히 보였고 \cite{bianchini2005inside}, Robust PageRank는 한 노드의 점수 중 소수의 기여자에게서 올 수 있는 몫을 제한한다 \cite{andersen2008robust}. \ref{sec:method}절은 이 생각들을 컨트랙트에 적용한다.
